\chapter{フレネ・セレの公式と応用}

\section{フレネ・セレの公式}

小節\ref{subsection:differential-of-frames}で示した公式を並べてみましょう。
\begin{align*}
  \frac{de_1(s)}{ds}&=\kappa(s)e_2(s)\\
  \frac{de_2(s)}{ds}&=-\kappa(s)e_1(s)+\tau(s)e_3(s)\\
  \frac{de_3(s)}{ds}&=-\tau(s)e_2(s)
\end{align*}
なんだか行列みたいですね。
行列みたいに書いてみましょう。
\[
  \frac{d}{ds}
  \begin{pmatrix}
    e_1(s)\\
    e_2(s)\\
    e_3(s)
  \end{pmatrix}
  =
  \begin{pmatrix}
    0 & \kappa(s) & 0 \\
    -\kappa(s) & 0 & \tau(s) \\
    0 & -\tau(s) & 0
  \end{pmatrix}
  \begin{pmatrix}
    e_1(s)\\
    e_2(s)\\
    e_3(s)
  \end{pmatrix}
\]
これは小節\ref{subsection:differential-of-frames}の計算結果を行列の形に並べたに過ぎないですが、曲線を文字通り特徴づける、意味ある等式です。
とても偉いので、名前を付けましょう。
\begin{definition}
  $I$を区間、$\gamma:I\to\mathbb{R}^3$を弧長パラメータ付け可能な$C^3$級曲線であって、弧長パラメータ$s=s(t)$に関して$d^2\gamma(t(s))/ds^2\neq0$を満たすとする。
  このとき直交行列
  \[
    \begin{pmatrix}
      e_1(s)\\
      e_2(s)\\
      e_3(s)
    \end{pmatrix}
  \]
  を\textbf{フレネ標構造}と呼ぶ。
\end{definition}
各$e_i(s)$は$\mathbb{R}^3$の元なので、フレネ標構は3次行列です。
また
\[
  e_i(s)\cdot e_j(s)=\begin{cases}
    1 & (i=j)\\
    0 & (i\neq j)
  \end{cases}
\]
なので、フレネ標構は直交行列なわけです。

フレネ標構に関して、次の定理が成り立つことは既に証明しました。
これは\textbf{フレネ・セレの公式}と呼ばれます。
\begin{theorem}[フレネ・セレの公式]
  $I$を区間、$\gamma:I\to\mathbb{R}^3$を弧長パラメータ付け可能な$C^3$級曲線であって、弧長パラメータ$s=s(t)$に関して常に$d^2\gamma(t(s))/ds^2\neq0$を満たすとする。
  また、
  \begin{align*}
    e_1(s)&:=\frac{d\gamma(t(s))}{ds}\\
    e_2(s)&:=\frac{d^2\gamma(t(s))/ds^2}{|d^2\gamma(t(s))/ds^2|}\\
    e_3(s)&:=e_1(s)\times e_2(s)
  \end{align*}
  をそれぞれ$\gamma$の単位接ベクトル、主法線ベクトル、従法線ベクトルとし、
  \begin{align*}
    \kappa(s)&:=\left|\frac{d^2\gamma(t(s))}{ds^2}\right|\\
    \tau(s)&:=e_3(s)\cdot\frac{de_2(s)}{ds}
  \end{align*}
  をそれぞれ$\gamma$の曲率、捩率とする。
  このとき、フレネ標構に関して
  \[
    \frac{d}{ds}
    \begin{pmatrix}
      e_1(s)\\
      e_2(s)\\
      e_3(s)
    \end{pmatrix}
    =
    \begin{pmatrix}
      0 & \kappa(s) & 0 \\
      -\kappa(s) & 0 & \tau(s) \\
      0 & -\tau(s) & 0
    \end{pmatrix}
    \begin{pmatrix}
      e_1(s)\\
      e_2(s)\\
      e_3(s)
    \end{pmatrix}
  \]
  が成り立つ。
\end{theorem}

\section{フレネ・セレの公式を解く}

「公式」を解くって何？って話ですが、フレネ・セレの公式を見ていると「先に曲率と捩率が与えられているとすれば、これは単なる微分方程式じゃね？」ってことです。
例えば形式的に
\begin{align*}
  E(s)&:=
  \begin{pmatrix}
    e_1(s)\\
    e_2(s)\\
    e_3(s)
  \end{pmatrix}\\
  M(s)&:=
  \begin{pmatrix}
    0 & \kappa(s) & 0 \\
    -\kappa(s) & 0 & \tau(s) \\
    0 & -\tau(s) & 0
  \end{pmatrix}\\
\end{align*}
とおけば、フレネ・セレの公式は非常に単純に見えるタイプの微分方程式
\[
  \frac{dE(s)}{ds}=M(s)E(s)
\]
に見えます。
もし$E(s)$や$M(s)$が単なる関数だったら
\begin{align*}
  &\frac{dE(s)}{ds}=M(s)E(s)\\
  \therefore&\frac{dE}{E}=M(s)ds\\
  \therefore&\int\frac{dE}{E}=\int M(s)ds\\
  \therefore&\log(E(s))=\int M(s)ds+C\\
  \therefore&E(s)=E_0\exp\left(\int M(s)ds\right)
\end{align*}
という風に解きたいところです。
しかし悲しいかな、$E(s)$も$M(s)$も行列ですから、もう少し慎重にやらねばなりません。
